\section{Regularization and restoration of the curvature invariant}\label[appendix]{sec:regularization}

The minimizing partition is smooth away from a singular set, whereas the
next cutting step needs smooth complete packets.  Shielding has made the
closed guard disjoint from the singular set; the first lemma below obtains
the regular tubes needed for comparison.  We construct globally smooth
partitions that accurately approximate the true partitions near these
guards.  We then transfer the scalar form to the new walls and joined
sides, smooth the cylinder seams, and convert the form bound into a
pointwise bound by adding circle factors.  Throughout this section a stage
$r$ and its current label $j_0$ are fixed.  The profiles and
cores of the remaining labels have already been updated by
\Cref{shield:global}.  Write
\[
\begin{gathered}
 \widetilde\Pi_j=\Pi_j+D_j,\qquad
 G'_j=E\setminus(U_j^{+,\mathrm{updated}}\cup U_j^{-,\mathrm{updated}}),\\
 H_E^*=\left\{x\in E:
       \#\{j\in J_r\setminus\{j_0\}:x\in G'_j\}\ge m-2\right\}.
\end{gathered}
\]
Here $m=\dim E$ is the parent base dimension.  The set $H_E^*$ is closed
and misses the singular set of the true boundary.  A nonpositive counting
threshold means the whole packet.  Scalar curvatures, mean curvatures,
Laplacians, and measures include the old torus factors; functions and tests
are invariant in those factors.  Constants in this section may depend on a
packet and its finite incidence diagram.  At the first, radial step the
corresponding constants are assigned on enlarged compact radial blocks.
In particular no uniform geometry bound for the original manifold is used.

Use adjacent parameters $0\le t_r<t_{r+1}\le1$, and put
\[
 \Delta_r=t_{r+1}-t_r,\qquad
 a_r=.60+.03t_r,\qquad b_r=1.12-.05t_r,\qquad
 \lambda_r=.25-t_r/8,
\]
\[
 \eta_r=\Delta_r/200,\qquad
 e_0=(\lambda_r-\lambda_{r+1})/100=\Delta_r/800,
 \qquad C=3.99.
\]
Only positivity of the particular increment is needed; none of the
geometric choices is required to be uniform in $r$.  The proofs below
are stated for finite label data.  Their constant-preserving form,
spelled out in \Cref{regular:constant-preservation}, will also be used
for a finite exceptional list together with a summable constant tail.

\subsection{Guarded partitions and transfer of the raw form}

\begin{lemma}[Uniform control on the guarded region]\label{regular:guards}
Fix a finite integer $J\ge4$.  On each parent packet $E$, let
$\varsigma_E$ be signed distance to the true boundary support, truncated
at $\pm1$, positive in the plus phase.  If the support is empty on a
component, use the corresponding constant sign.  After the radial split,
the following choices can be made on each whole parent packet.
\begin{enumerate}[label=(\roman*)]
\item There is a positive $\rho_E<1/10$ such that near
\[
 \{x:\dist(x,H_E^*)<4\rho_E,\ |\varsigma_E(x)|<4\rho_E\}
\]
the true boundary is regular, has an embedded normal tube, and signed
distance has bounded derivatives through order $J+2$ on a slightly
smaller tube.  The current forced core closures have positive uniform
distance from the true boundary.
\item On an elementary collar $B\times(T,\infty)$, signed distances
converge uniformly to the pullback of $\varsigma_B$ as $T\to\infty$.
They converge in $C^J$ on the guarded tubes, with the coface tube chosen
inside the face's tube.  These statements are uniform over the whole
face, including its noncompact ends.
\end{enumerate}
For the initial radial parent, assertion \textup{(i)} holds separately
on each enlarged compact radial block, with its own radius and derivative
bounds; assertion \textup{(ii)} is not required.
The radii depend on the true partitions and the updated guards, and can
be chosen before any smooth-approximation tolerances.
\end{lemma}

\begin{proof}
On a compact set, closedness of $H_E^*$ and its disjointness from the
singular set give a regular neighborhood of every point where the guard
meets the boundary.  A finite subcover and the normal inverse function
theorem give a positive tube radius and all the asserted finite-order
bounds there.  The same compactness argument applies to the separation
of the current core closures, which lie inside the barrier domains of
\Cref{cut:pressure,cut:gmt} and whose neighborhoods miss the support.

We explain why these statements remain uniform at infinity.  If a
sequence escapes in $E$, the finite face description in
\Cref{packet:escape} gives, after passage to a subsequence, a product
window over a further face $F$ in which the $F$ coordinates converge
in a compact set and all the remaining normal coordinates tend to
infinity.  Bounded metric balls lie in these product windows eventually.
The true phases converge there to the product of the true phase of
$F$, by \Cref{cut:end-convergence}.  The updated cores and their
complements are exact products.  Consequently a limit of guarded points
lies in $H_F^*$: a count of at least $m-2$ is at least
$\dim F-2$.  The limiting boundary point is therefore regular.
Smooth, single-sheet convergence near this point supplies an embedded
normal tube of a fixed positive radius, and excludes additional support
in a smaller neighborhood.  In particular no sequence at distance
tending to zero from both the guard and the support can have degenerating
tube radius or unbounded prescribed finite derivatives.  This proves
the first assertion by contradiction, with buffers reduced if necessary.
The barrier statement at infinity follows in the same way from its
face statement and phase/support convergence.

For uniform convergence in an entire elementary collar, suppose a
discrepancy persists at $(x_i,t_i)$ with $t_i\to\infty$.  If $x_i$
stays in a compact part of $B$, local product convergence gives a
contradiction.  Otherwise descend through the finite face diagram of
$B$.  Commutation of its collars with the $t$ collar puts both the
partition of $E$ and that of $B$ over the same further face $F$,
with compact tangential coordinates and all escaping coordinates tending
to infinity.  Both limiting partitions are the same unique product
partition.  Local support convergence implies convergence of truncated
distances: a nearest point at distance less than $1$ is detected in a
slightly larger fixed ball, and if there is no such point both truncated
distances approach the saturated value.  Away from the limiting support,
phase convergence fixes the sign.  This proves uniform $C^0$
convergence.  On a guarded tube the limiting boundary is regular; smooth
graph convergence and the normal inverse function theorem give $C^J$
convergence of signed distance on smaller tubes.

Choose radii from the deepest faces upward.  Each coface radius is
smaller than a fixed fraction of every relevant further-face radius.
For example, a point within $\rho_E$ of $H_E^*$ projects, far out, to
within $2\rho_E$ of $H_B^*$, and $|\varsigma_E|<\rho_E$ implies
$|\varsigma_B|<2\rho_E$ there.  These are statements about the fixed true
geometry and exact cores, without approximation tolerances.  This order
also handles a simultaneous escape in several coordinates and proves the
last assertion.  At the radial step the same proof on compact blocks
requires no assertion about an initial spatial end.
\end{proof}

\begin{lemma}[Compatible smooth partitions]\label{regular:partition}
There are compatible smooth defining functions $g_E$ with $0$ a regular
value, positive and negative on the respective current forced core
closures, and with $L_E=g_E^{-1}(0)\subset\operatorname{int}G_{j_0}$.
After the radial split they can be chosen uniformly as close as
prescribed to $\varsigma_E$ in value and in $C^J$ on smaller guarded
tubes.  In a buffered neighborhood of $H_E^*$, a smooth diffeomorphism
arbitrarily close to the identity in $C^{J-1}$ identifies the smooth
partition with the true partition.  It uses the same map on the wall
and on the boundaries of its two sides.  At the initial radial step,
the approximation tolerances and comparison buffers are chosen separately
on enlarged compact radial blocks.  The defining functions, but not this
comparison map, are exact products in sufficiently far old collars.
\end{lemma}

\begin{proof}
First choose all the geometric radii in \Cref{regular:guards}.  Prescribe
positive value and $C^J$ error targets on the guarded regions, small
enough also to preserve the core signs.  Allocate smaller targets to
faces to meet the requests of all their cofaces, including through
chains, as in \Cref{packet:approximation}.  There are only finitely many
such requests per face and finite depth.  This allocation does not
change any tube radius.

Work now from deeper faces upward.  Smooth $\varsigma_E$ by local
convolution and a smooth partition of unity, with arbitrarily small
value error everywhere and derivative error on the regular buffers.
These errors can be prescribed separately: use the smooth function
itself on an inner regular buffer and blend to a value approximation
outside a larger one.  On a far collar blend to the previously chosen
face function.  \Cref{regular:guards} makes the discrepancy from the
face's true signed distance arbitrarily small by moving the collar
start.  The discrepancy of the already chosen face function from that
distance is instead paid from its preallocated error target.  These
two errors are kept separate.

More explicitly, take cutoffs $f_i$ from $0$ to $1$ in the elementary
normal coordinates, in valid far product ranges.  Use weights
\[
 w_0=\prod_i(1-f_i),\qquad
 w_i=f_i\prod_{a<i}(1-f_a).
\]
They sum to $1$.  All derivatives through order $J$ are bounded by
constants fixed before the targets are allocated (translated cutoffs
have the same bounds).  Thus blending costs a fixed multiple of the
participating errors.  Face copies agree on all relevant overlaps once
the starts exceed their product thresholds.  When any $f_i=1$, every
participating copy equals that face copy on the overlap; hence the
result is exactly compatible.  This proves the claimed approximations.

Near a possible zero in a guarded tube, the defining function has
nonzero normal derivative and its zero is a small normal graph over the
true wall.  Off the tube, uniform value closeness preserves the phase
where its sign is required.  The copied end functions already have
regular zeros.  On the remaining compact region, a relative regular-value
perturbation makes all zeros regular without changing the copied ends.
For completeness, if $g$ is the preliminary function, choose a smooth
$\chi\ge0$, positive on its nonregular zeros and zero near the regions
to be kept fixed.  For an arbitrarily small regular value $c$ of
$-g/\chi$ on $\{\chi>0\}$, the function $g+c\chi$ has regular zeros
there.  Where $\chi=0$, its differential is that of $g$, since
$d\chi=0$ there.  Bounded derivatives of $\chi$ and small $c$ preserve
all strict error and sign requirements.  On the initial space use a
positive smooth $\chi$ with locally prescribed small derivatives and
the same regular-value argument; the radial separators are preserved.

If the perturbed wall is $t=\alpha(x)$ in a true normal tube, the map
\[
 (x,t)\longmapsto (x,t-\chi_1(x,t)\alpha(x))
\]
with $\chi_1=1$ on the needed graph region and supported inside a larger
regular buffer is a diffeomorphism for sufficiently small graph norm.
It carries the graph to $t=0$, preserves its two local sides, and tends
to the identity in $C^{J-1}$.  Choose the buffer cutoff on the regular
wall slightly beyond the projection of the smaller guard neighborhood
and extend by the identity.  Uniform regular charts give bounded
cutoff derivatives.  The same formula restricts to the wall used as a
cylinder base.  No product property of this auxiliary comparison is
asserted or needed.
\end{proof}

Perform the split of \Cref{packet:split} using these defining functions.
For the moment use the raw metric: the old metric on the bulk and the
frozen induced metric on each fresh infinite cylinder.  Let $p_A$ denote
restriction/inclusion on a wall or bulk and cylinder projection on a
fresh cylinder.  Pull back the remaining profiles, cores, and
$\widetilde\Pi_j$ by $p_A$.  These data are continuous at a seam and
compatible in old ends.  Define the new required set $D_A$ from the
remaining cores and the new base dimension.
The raw forms of \Cref{cut:raw-domination} were defined on the true
minimizing wall and its joined sides.  We first estimate those forms,
then use the comparison map to transfer the estimate to the raw metrics
of the smooth split just constructed.

\begin{lemma}[Counts and transfer of the raw form]\label{regular:transfer}
The defining-function accuracies can be chosen so that on a uniform
neighborhood of $D_A$ in each new packet the raw scalar form controls
\begin{equation}\label{regular:raw-target}
 \lambda_{r+1}+(1+3\eta_r)b_{r+1}\sum_{j\ne j_0}\widetilde\Pi_j,
\end{equation}
with a strictly positive mass and gradient reserve.  The form has
gradient coefficient $C=3.99$ and includes the seam term
$2\int H_{\rm out}u(0)^2$ on a joined side.  Tests are compactly
supported, with matching seam traces.
\end{lemma}

\begin{proof}
For a side packet $\dim A=m$, a point in $D_A$ projects to at least
$m-1$ remaining undecided labels, hence into $H_E^*$ and the parent's
old required set.  For a hit packet $\dim A=m-1$, it projects to at
least $m-2$ remaining undecided labels.  The wall lies in $G_{j_0}$,
which supplies the additional old undecided label.  The projection is
again in both required sets.  These inclusions are unchanged on the
fresh cylinder, where all relevant data are pulled back from its base.
All these maps decrease lengths in the raw metrics; therefore a small
neighborhood of $D_A$ projects into small neighborhoods of the two
parent sets.  The closed-set count, rather than a continuity claim for
the counting function, is what is used here.

On compared true pieces in a sufficiently small neighborhood of the
old required set, continuity and the bounded derivatives of the fixed
parent data give
\[
 R\ge\lambda_r+b_r\sum_j\Pi_j-e_0,
 \qquad e_0=\frac{\lambda_r-\lambda_{r+1}}{100}.
\]
By \Cref{cut:raw-domination,cut:reserve}, the raw form on these pieces
dominates $R-1.02\Pi_{j_0}$, with spare gradient coefficients $.94$
on bulk and $1.99$ on a wall or cylinder.  The bound in
\Cref{shield:global} pays for the additional potential
$1.16\sum_{j\ne j_0}D_j$ from these spare gradient terms and a mass
loss of at most $.9e_0$.  Since
$b_r>1.02$, dropping the remaining nonnegative current-label term
leaves
\begin{align}
 Q_{\rm true}(u)&\ge
 \int\left(\lambda_r-2e_0+
       b_r\sum_{j\ne j_0}\Pi_j+
       1.16\sum_{j\ne j_0}D_j\right)u^2
       +.2\int|\nabla u|^2.\label{regular:true-budget}
\end{align}
Indeed the unrounded gradient remainders are
$.94-.9\cdot.8=.22$ and $1.99-.9\cdot1.7=.46$, and the mass loss is
at most $1.9e_0$.  The exact numerical margins are
\begin{equation}\label{regular:eta}
\begin{aligned}
 a_{r+1}-(1+\eta_r)a_r
   &=\frac{3\Delta_r(180-t_r)}{20000}
     \ge .02685\Delta_r>0,\\
 b_r-(1+3\eta_r)b_{r+1}
   &=\frac{\Delta_r[664+15(t_r+\Delta_r)]}{20000}
     \ge .0332\Delta_r>0,\\
 1.16-(1+3\eta_r)b_{r+1}
   &=\frac{800+1000t_r+664\Delta_r+
                  15\Delta_r(t_r+\Delta_r)}{20000}>0.
\end{aligned}
\end{equation}
These identities use $0\le t_r<t_r+\Delta_r\le1$.  Thus
\Cref{regular:true-budget} exceeds
\Cref{regular:raw-target} by at least $98e_0$ in mass and $.2$ in
gradient, in addition to strictly positive potential margins.

Use the comparison diffeomorphism of \Cref{regular:partition} on bulk
and on the common wall/cylinder base, leaving the cylinder coordinate
and torus variables unchanged.  A test supported in a smaller buffered
neighborhood of $D_A$ becomes a compact matching test on the true
regular pieces, extended by zero outside the comparison region.  For
a test in one component set it to zero on other components.  In the
bulk away from the wall the map is the identity and the phase signs
already agree.  All supports stay away from singular points and their
cylinders.

The pulled-back metric, volume weight, potentials, intrinsic stabilized
wall curvature, and outward mean curvature converge uniformly with the
chosen finite derivatives.  For example the induced scalar curvature
can be computed from Gauss' equation, so the second fundamental form
and the ambient curvature give the requisite bounds.  Consequently,
for arbitrarily small $\delta>0$ the discrepancy of the compared forms
is bounded by
\[
 \delta\left[\int(u^2+|\nabla u|^2)+\int_{\rm seam}u(0)^2\right].
\]
The last term is controlled using the first unit interval of the raw
product cylinder:
\begin{equation}\label{regular:trace}
 \int_{\rm seam}u(0)^2
 \le 2\int_{S\times[0,1]}(u^2+|\partial_tu|^2).
\end{equation}
This follows by writing $u(0)=u(t)-\int_0^t\partial_tu$ and integrating
in $t$.  It needs no zero boundary condition.  Choosing all approximation
errors sufficiently small leaves, for example, a reserve of $90e_0$
in mass and $.15$ in gradient.  These choices can be prescribed on
parent buffers before the individual new components are known: the
projection inclusions place every relevant point there.  At the radial
step the enlarged compact blocks supply the bounds; preserved
separators trap each new component in one such block before cylinder
attachment.  Finally impose all coface requests through the finite
tolerance allocation already used in \Cref{regular:partition}.
\end{proof}

\subsection{Seam smoothing and restriction of profiles}

The proof below uses scalar-curvature concentration at a mean-curvature
jump, as in the corner smoothing of
\citet[Sections~2--3]{miao2002Corners}.  Here the smoothing width varies
tangentially, and the form estimate must preserve compatible product ends.

\begin{lemma}[Smoothing a seam with variable width]\label{regular:seam}
The raw joined metrics admit smooth compatible metrics with arbitrarily
small $C^0$ error, and their scalar forms converge to the raw forms in
the following quantitative sense.  On each joined packet and for every
$\delta>0$, the smoothing can be chosen so that the absolute form
error is at most
\[
 \delta\int(u^2+|\nabla u|^2)
\]
for every compact matching test, using the raw metric to measure the
right side.  The old warped-product structure and every sufficiently
far old product end are preserved.  The choices on a wall can satisfy
the requests of both incident sides at once.
\end{lemma}

\begin{proof}
In a normal collar of a smooth wall write the old stabilized metric
as $ds^2+\beta_s$ for $s\le0$, where $s$ increases outward.  Here
$\beta_s$ includes the old torus factors.  The raw joined metric is
$ds^2+\beta_{\min(s,0)}$.  Choose a smooth concave function $\Phi$
equal to $t$ for $t\le-1$ and to $0$ for $t\ge1$, with
$0\le\Phi'\le1$.  Such a function is obtained by integrating a smooth
decreasing step function symmetric about $0$.  It satisfies
$\Phi\le0$ and $\int\Phi''=-1= -\int|\Phi''|$.

For a smooth positive function $\epsilon(x)$ on the wall put
\[
 \phi(x,s)=\epsilon(x)\Phi(s/\epsilon(x)),\qquad
 G_\epsilon=ds^2+\beta_{\phi(x,s)}.
\]
It equals the raw metric outside $|s|<\epsilon(x)$ and differs from it
in $C^0$ by $O(\sup\epsilon)$.  In a fixed tangential chart the extra
derivatives of $\phi$ satisfy
\begin{align*}
 |\partial_x\phi|&\le c|d\epsilon|,\\
 |\partial_x^2\phi|&\le
 c\bigl(|d^2\epsilon|+|d\epsilon|^2/\epsilon\bigr).
\end{align*}
Choose $\epsilon=\zeta^2$ with $\zeta>0$ as small in $C^2$ as
prescribed.  Then $|d\epsilon|^2/\epsilon=4|d\zeta|^2$ and
$d^2\epsilon=2\zeta d^2\zeta+2d\zeta\otimes d\zeta$.
Thus these tangential errors tend to zero.  On a packet, positive
compatible choices have a positive infimum and bounded derivatives by
the compact-part/product-end description.  Their size, but not the
resulting higher metric bounds, can be made arbitrarily small.  Normal
collar widths are also positive packetwise by the same description.

Let $B^{\rm old}_\tau,H^{\rm old}_\tau$ be the slice second
fundamental form and mean curvature of $ds^2+\beta_s$.  Exactly,
\[
 B_s=\phi_sB^{\rm old}_{\phi},\qquad
 H_s=\phi_sH^{\rm old}_{\phi}.
\]
For a metric $ds^2+\beta(x,s)$ the Gauss and normal variation
identities give
\[
 R(G_\epsilon)=R(\beta(x,s))-|B_s|^2-H_s^2-2\partial_sH_s.
\]
The tangential derivative bounds just proved show that all terms on
the strip are uniformly bounded except
\begin{equation}\label{regular:spike}
 -2\phi_{ss}H^{\rm old}_{\phi}.
\end{equation}
In particular the formula differentiates $H_s$ only in $s$;
potentially large mixed $xs$ derivatives do not give another scalar
spike.  The remaining contribution from this derivative is
$-2\phi_s^2\partial_\tau H^{\rm old}_\phi$, which is bounded.

Let $J_\epsilon(x,s)$ be the density of $\beta_\phi$ relative to
the zero slice.  Uniformly on the strip,
$J_\epsilon=1+O(\epsilon)$ and
$H^{\rm old}_\phi=H^{\rm old}_0+O(\epsilon)$.
For each fixed $x$,
\[
 \int-\phi_{ss}\,ds=1,\qquad
 \int|\phi_{ss}|\,ds=1.
\]
Writing $E_x=\int_{-\epsilon(x)}^{\epsilon(x)}|\partial_su|^2ds$,
the fundamental theorem of calculus gives
\[
 |u(s)-u(0)|^2\le\epsilon(x)E_x,
 \quad
 |u(s)^2-u(0)^2|
 \le2\sqrt{\epsilon(x)}|u(0)|E_x^{1/2}+\epsilon(x)E_x.
\]
After integrating against $|\phi_{ss}|ds$, the error in replacing
$u(s)^2$ by $u(0)^2$ is at most
$c\sqrt{\sup\epsilon}\,(u(0)^2+E_x)$, for $\sup\epsilon\le1$.
The density and mean-curvature errors cost at most
$c\sup\epsilon\,(u(0)^2+E_x)$.  Thus the form contribution of
\Cref{regular:spike} converges to
$2\int H^{\rm old}_0u(0)^2$, with the outward sign required in
\Cref{cut:side-estimate}.

The bounded strip terms have vanishing form cost because
\[
 \int_{-\epsilon(x)}^{\epsilon(x)}u(s)^2ds
 \le c\epsilon(x)u(0)^2+c\epsilon(x)^2E_x.
\]
The metric and density errors in the gradient and fixed bounded
potential terms tend to zero in $H^1$ form norm.  Integrate these
estimates in $x$ and use \Cref{regular:trace}; the stated error bound
follows.  One first proves it for smooth matching tests and then for
the matching $H^1$ completion by density.

Choose $\zeta$ on new hit packets, from deeper faces upward, after
propagating the two incident-side requests through their finite
diagrams.  The function is an exact pullback in old product ends, so
the displayed smoothing is also an exact product there.  It is
unchanged far out in the fresh cylinder coordinate.  Since the normal
coordinate and old metric already split off every old normal factor,
the operation commutes with all the inherited collar identifications.
\end{proof}


\begin{lemma}[Approximation that reproduces constants]\label{regular:constant-preservation}
A continuous scalar function on a packet can be approximated by smooth
functions, with any prescribed positive value error, by operators that
reproduce constants.  For compatible input on a packet diagram, the
approximations can be chosen compatible in all sufficiently far product
ends of that diagram.  If the input and all face copies
used at a packet equal the same constant $c$, the result equals $c$
there exactly.  The same assertion holds componentwise for a finite list
of functions and for a list having only finitely many nonconstant entries
on each prescribed finite directed face/coface diagram.
\end{lemma}

\begin{proof}
Choose a locally finite coordinate cover, relatively compact smaller
charts, and a smooth partition of unity $\{\psi_\alpha\}$ subordinate
to the smaller charts.  Extend each chart slightly and average the
function against a smooth nonnegative kernel of integral one, at a
sufficiently small chart-dependent scale.  Call the resulting local
function $T_\alpha f$.  Uniform continuity on the chart closure makes
$|T_\alpha f-f|$ smaller than the prescribed error on
$\supp\psi_\alpha$.  Thus
\[
                    Tf=\sum_\alpha\psi_\alpha T_\alpha f
\]
is a smooth approximation with that error.  Every local average sends
$c$ to $c$, and $\sum\psi_\alpha=1$, so $Tc=c$ exactly.  Locally
prescribed derivative accuracy is obtained in the same way wherever
$f$ is already smooth, using a slightly larger regular buffer.

Now use the finite tolerance allocation of
\Cref{packet:approximation}, and blend from deeper faces upward with
the weights $w_i$ in the proof of \Cref{regular:partition}.  Their sum
is one.  Consequently blending preserves a common constant as well as
exact face identities.  The finite exceptional union over the prescribed
directed dependencies is chosen before applying this construction.
For every other entry the entire local input and all participating face
copies coincide with the same constant, so all its averaging and
blending errors are zero.  No minimum of tolerances over these constant
entries is needed.  On the initial radial space the construction is
local on buffered compact blocks and uses a locally finite chart cover.
It has the same constant-reproduction property there.
\end{proof}

On a new packet $A$, write $R_{\rm old}$ and $\lap_{\rm old}$ for the
scalar curvature and invariant Laplacian of the smoothed split metric
with its retained torus factors, before any new circle factors are added.

\begin{lemma}[Restricted profiles and smooth potentials]\label{regular:profiles}
After the seam smoothing the pulled-back profiles and cores satisfy
the stage $r+1$ profile conditions with smooth compatible potentials
$\Pi_j^{\rm new}$ chosen so that
\[
 (1+\eta_r)\widetilde\Pi_j
 \le\Pi_j^{\rm new}\le(1+2\eta_r)\widetilde\Pi_j.
\]
If
\begin{equation}\label{regular:Vstar}
 V_A^*=\lambda_{r+1}+b_{r+1}\sum_{j\ne j_0}\Pi_j^{\rm new},\qquad
 L_A^0=-C\lap_{\rm old}+R_{\rm old}-V_A^*,
\end{equation}
then on some positive uniform neighborhood of $D_A$ its form is bounded
below by a positive multiple of $\|u\|_2^2$.
\end{lemma}

\begin{proof}
We justify restriction to a hypersurface, which does not follow merely
by restricting an ambient almost-everywhere inequality.  If $f$ is
locally Lipschitz on an open Riemannian set and $|df|\le F$ almost
everywhere for continuous $F$, then at every point its local Lipschitz
slope is at most $F$.  In a small coordinate ball, convolve $f$ on a
slightly larger ball.  The derivative bound of the convolutions is at
most the supremum of $F$ there times a metric distortion factor tending
to $1$ as the ball shrinks.  Integrating this bound along straight
segments, passing to the uniform limit, and shrinking the ball proves
the assertion.  It follows also that every rectifiable curve satisfies
the corresponding integrated bound, by subdividing into such balls.

Apply this on each open finite band of a profile, with
$F=(a_r\mu_j^2+\widetilde\Pi_j)/2$.  Restriction to the wall,
inclusion of bulk, and projection on a cylinder therefore preserve the
bound everywhere in local slope, and hence almost everywhere for the
restricted locally Lipschitz function.  Finite truncations remain
locally Lipschitz, including where a pole is present.  Across a seam
the same conclusion follows by splitting a path into its bulk and
cylinder portions; traces agree.  A sufficiently small $C^0$ metric
change increases all these slopes by at most $1+\eta_r$.

For each label $j$, the raw potential has its own positive lower
bound $p_j^*>0$.  Apply \Cref{regular:constant-preservation} to
$(1+3\eta_r/2)\widetilde\Pi_j$, with strict value error smaller than
$\eta_r p_j^*/2$.  This proves the displayed relative bounds.  Different
labels may require different accuracies; a common positive lower bound
for an infinite label list is unnecessary.  When the raw potential is
a spatial constant $c_j$ on the whole relevant incident diagram, choose
its new value to be exactly $(1+3\eta_r/2)c_j$.  By
\Cref{regular:eta},
\[
 2|\nabla\mu_j|_{\rm new}
 \le (1+\eta_r)(a_r\mu_j^2+\widetilde\Pi_j)
 \le a_{r+1}\mu_j^2+\Pi_j^{\rm new}.
\]
Continuous projections preserve the inclusion of every core closure
inside its pole interior and preserve all prescribed $q$ constraints.

Choose the seam form error small relative to the reserves of
\Cref{regular:transfer}.  The potential in \Cref{regular:Vstar} is at
most that in \Cref{regular:raw-target}; the extra $\eta_r$ supplies
additional margin for the potential smoothing.  Metric equivalence
allows the required-set neighborhood to be slightly reduced after
smoothing.  Thus one may retain, for example, an $e_0\|u\|_2^2$
lower bound for the form of $L_A^0$ on tests in that neighborhood.
No compatibility of the neighborhoods themselves is asserted.
\end{proof}

We have obtained a smooth packet metric and the next profile
inequalities. The scalar lower bound is still only a quadratic-form
bound near the new required set. We next extend its positive gap to
the whole packet by a correction that vanishes near that set, solve a
bounded elliptic equation, and approximate its solutions compatibly.
The final group of $399$ circle factors converts this form information
into the pointwise scalar inequality.

\subsection{Compatible corrections and positive supersolutions}

At this point the smooth split packets have compact truncations and
exact product ends by \Cref{packet:split}.  At the radial step the
preserved separating bands give the compact initial pieces.  Hence
\Cref{packet:geometry} already applies to these metrics and retained
warps, before the following correction is constructed.  In particular
all bounds used below hold separately on each whole packet; they are
not assumptions on the original manifold.

\begin{lemma}[A compatible correction with a global gap]\label{regular:correction}
There are smooth nonnegative bounded compatible potentials $K_A$,
with all derivatives bounded packetwise, such that $K_A=0$ in a
positive uniform neighborhood of $D_A$ and
\begin{equation}\label{regular:global-gap}
 Q_A(u):=\int\bigl(C|\nabla u|^2+
       (R_{\rm old}-V_A^*+K_A)u^2\bigr)
 \ge g_A\|u\|_2^2\qquad(g_A>0)
\end{equation}
for every compact smooth test.
Let $\omega_A^{\rm old}$ be the product of the retained circle length
functions for the smoothed metric on $A$, before the new circle
factors at this split are appended.  With $K_A$ acting by
multiplication, define
\[
 L_A=L_A^0+K_A
     =-C\lap_{\rm old}+R_{\rm old}-V_A^*+K_A.
\]
Its quadratic form in the measure
$\omega_A^{\rm old}\,\dd\vol_A$ is $Q_A$.
In an exact end over $B$, the base metric is a product and every
retained warp pulls back from $B$.  Consequently
\[
 \lap_{{\rm old},A}=\lap_{{\rm old},B}+\lap_{\rm normal},
 \qquad L_A=L_B-C\lap_{\rm normal}.
\]
\end{lemma}

\begin{proof}
First disregard compatibility on a fixed packet.  Put
$W_A=R_{\rm old}-V_A^*$.  Choose a smooth angle $\theta$ equal to
$0$ near $D_A$ and $\pi/2$ outside the neighborhood of local positivity,
with bounded derivatives and with the support of $\cos\theta$ inside
that neighborhood.  Uniform charts and the positive neighborhood
width give such a function by smoothing a distance cutoff with
buffers.  Set
\[
 K_{\rm ind}=J\sin^2\theta+C|\nabla\theta|^2,
 \qquad J>\max\{0,-\inf W_A\}.
\]
The identity for the squared partition $\cos\theta,\sin\theta$
shows that $Q_A^0+K_{\rm ind}$ has gap
$\gamma_A=\min\{e_0,J+\inf W_A\}>0$.  Indeed the gradient
partition error is exactly $C|\nabla\theta|^2u^2$; local positivity
controls the cosine part and $W_A+J>0$ controls the sine part.
The correction is zero on a smaller uniform neighborhood of $D_A$.

Proceed from deeper faces upward.  In sufficiently far elementary
collars the already constructed $K_B$ copies agree on intersections:
on a common quadrant each is the pullback of the same further-face
potential.  They consequently define a \emph{single} nonnegative
function $K_{\rm end}$ on the collar union.  Copies are not added
together on an intersection.  The metric, measure, and $V_A^*$ are
likewise exact products there.  In particular $V^*$ has the same
formula on a face: $C$, $\lambda_{r+1}$, $b_{r+1}$, the label set,
and the old torus list do not depend on base dimension, after padding
unused flat factors on each finite compared diagram.

For the finitely many elementary collars choose angles $\theta_i$
which rise from $0$ to $\pi/2$ over a normal interval of length $L$,
inside the valid common product ranges.  Use the squared partition
\[
 \eta_0=\prod_i\cos\theta_i,\qquad
 \eta_i=\sin\theta_i\prod_{a<i}\cos\theta_a.
\]
It has $\sum_i\eta_i^2=1$ and
$\sum_i|d\eta_i|^2\le\sum_i|d\theta_i|^2\le c_A L^{-2}$;
the inequality follows successively from the two-term sine/cosine
identity, so no number of intersections is being summed separately.
The support of $\eta_0$ lies in a compact collar truncation.

Choose smooth multipliers $\alpha,\beta\in[0,1]$ with $\alpha=1$
on $\supp\eta_0$ and $\beta=1$ on the union of the face-cutoff
supports.  Let $\alpha$ vanish only farther out than those transition
intervals.  Let $\beta$ vanish on the inner side of a buffered collar
union, so $\beta K_{\rm end}$ extends smoothly by zero.  Both can be
formed from products of one-coordinate cutoffs.  Define
\[
 K_A=\alpha K_{\rm ind}+\beta K_{\rm end}.
\]
On the central support this dominates $K_{\rm ind}$; on each collar
support it dominates the single face copy.  Eventually
$\alpha=0,\beta=1$, giving exact compatibility.  Nonnegativity permits
the harmless overlap of these two terms in the finite transition
region.

This correction vanishes in a uniform neighborhood of $D_A$.  The
individual correction already does.  For the end term, a point of
$D_A$ in a product collar projects into $D_B$, because a face has a
smaller required counting threshold.  Start the support of $\beta$
with a buffer from the collar entry.  A point sufficiently close to
$D_A$ in this support then has a nearby required point in the same
collar, and their face projections are just as close.  It lies in the
vanishing neighborhood of $K_B$.  Take the minimum of the finitely
many face widths and entry buffers; at corners the copies agree, so
the same conclusion applies.

Apply the individual gap to $\eta_0u$.  Apply each face gap by slicing
to $\eta_iu$; its support is away from the collar entry, and the
normal gradient term is nonnegative.  If
$g_* =\min(\gamma_A,\{g_B\})>0$, the finite IMS identity gives
\[
 Q_A(u)\ge(g_*-Cc_A L^{-2})\|u\|_2^2.
\]
Choose $L$ so that the error is at most $g_*/2$.  All quantities used
for this choice are fixed before solving any elliptic equation.
Boundedness and the derivative assertions follow from compact parts
and finitely many product copies.  Empty required sets or absent faces
cause no change: one may take $\theta=\pi/2$ or omit the corresponding
cutoff.  This completes the compatible induction and proves the
operator identity.
\end{proof}

\begin{lemma}[The bounded gapped resolvent]\label{regular:resolvent}
Let $A$ be one of the resulting complete packets, with its weighted
base measure.  Suppose $W$ is smooth and bounded with bounded
derivatives, and $L=-C\lap_\omega+W$ has form gap $g>0$.  Then
there is a unique bounded smooth solution $s$ of $Ls=1$.
It satisfies $0<\inf s\le\sup s<\infty$ and has bounded derivatives
of every fixed order.  Bounded smooth supersolutions of $Lw\ge0$ are
nonnegative.  The same conclusions hold on Euclidean products of $A$.
\end{lemma}

\begin{proof}
Use the bounded geometry, positive bounded weight, and center-uniform
polynomial volume estimate of \Cref{packet:geometry}.  The form gap
and boundedness of $W$ make the closed form coercive on $H^1$:
combine $Q\ge g\|u\|_2^2$ with
$Q\ge C\|\nabla u\|_2^2-\|W_-\|_\infty\|u\|_2^2$.
It therefore has a self-adjoint inverse on $L^2$, of norm at most
$g^{-1}$.  If $f\ge0$, testing $Lu=f$ with the negative part gives
$Q(u_-)\le0$, so the inverse is positive.

Here is the decay estimate needed for bounded sources.  We use an
Agmon-type exponential-weight argument \citep{agmon1982}, proved here
for the weighted packet operator.  If $f$ is
supported in a metric unit ball $I$ and $u=L^{-1}f$, put
$\chi=\exp(\delta\min\{d(\cdot,I),D\})$, where $C\delta^2<g$.
The form identity for a bounded Lipschitz multiplier gives
\[
 Q(\chi u)=\langle Lu,\chi^2u\rangle+
                  C\int|\nabla\chi|^2u^2.
\]
It holds by compact approximation in $H^1$; on the support of $f$,
$\chi=1$.  Hence
\begin{equation}\label{regular:off-diagonal}
 (g-C\delta^2)\|\chi u\|_2^2
 \le\langle f,u\rangle\le g^{-1}\|f\|_2^2.
\end{equation}
Let $D\to\infty$.  The polynomial volume bound, uniform in the
center of $I$, implies
\[
 \int_A e^{-2\delta d(x,I)}\,\dd\mu(x)\le C_{A,\delta}<\infty.
\]
Indeed sum the exponential bound on successive unit annuli against
their polynomial volume upper bound.

Choose smooth compactly supported $0\le f_R\le1$ equal to $1$ on
an increasing compact exhaustion, and put $u_R=L^{-1}f_R\ge0$.
These are initially $H^1$ weak solutions.  For each fixed $R$, the
coefficients and $f_R$ are smooth, and weak interior regularity makes
$u_R$ smooth
\citep[Lecture~6, Theorem~1 and Corollary~1]{simonPDE}.
This assertion requires no derivative bounds for $f_R$ uniform in $R$.

For any $h\in L^2(I)$, self-adjointness,
\Cref{regular:off-diagonal}, and Cauchy--Schwarz with opposite
exponential weights give
\[
 |\langle u_R,h\rangle|
 =|\langle f_R,L^{-1}h\rangle|
 \le\frac{C_{A,\delta}^{1/2}}{\sqrt{g(g-C\delta^2)}}\|h\|_2.
\]
Thus $u_R$ has a uniform local $L^2$ bound on every unit ball, with
constant independent of its center and of $R$.  Choose the uniform
chart radius small enough that each larger chart ball lies in a
metric unit ball.  Writing $B_1$ and $B_{1/2}$ for the fixed rescaling
of two such nested chart balls, interior local boundedness gives
\[
 \|u_R\|_{L^\infty(B_{1/2})}
 \le c_A\bigl(\|u_R\|_{L^2(B_1)}
                      +\|f_R\|_{L^\infty(B_1)}\bigr).
\]
Here we use \citet[Lecture~18, Theorem~2, with $p=2$]{simonPDE};
the fixed radius and coordinate changes are absorbed in $c_A$.
The principal coefficients are
uniformly elliptic and bounded, the lower-order coefficients are
bounded, and the weight is bounded above and below.  Since
$0\le f_R\le1$, this gives a bound independent of the center and $R$.

Fix two nested relatively compact coordinate domains.  For all
sufficiently large $R$, the source is exactly $1$ on the larger domain.
The interior Schauder estimate
\citep[Lecture~12, Theorem~2]{simonPDE}, followed by differentiation,
therefore bounds every derivative of $u_R$ on the smaller domain,
uniformly for these $R$.  A diagonal subsequence converges smoothly
on compact subsets to a bounded nonnegative solution $s$ of $Ls=1$.
Applying the same estimates to this limiting equation in the uniform
packet charts gives bounded derivatives of every fixed order on $A$.

The maximum argument must allow $W$ to have either sign.  Let $w$ be
a bounded smooth supersolution, put $w_-=\max\{-w,0\}$, and fix
$0<\delta<\sqrt{g/C}$.  Write
$\psi=e^{-\delta d(x,x_0)}$ and take compact Lipschitz cutoffs
$\zeta_R$ equal to $1$ on $B(x_0,R)$, supported in $B(x_0,2R)$,
with $|d\zeta_R|\le 2/R$.  Testing the supersolution with
$(\psi\zeta_R)^2w_-$, and expanding the gradient, gives directly
\[
 g\|\psi\zeta_R w_-\|_2^2
 \le Q(\psi\zeta_R w_-)
 \le C\int |d(\psi\zeta_R)|^2w_-^2.
\]
All tests are compact, so no prior global gradient bound on $w$ is
needed.  Boundedness of $w$ and polynomial volume growth make
$\psi w_-$ square integrable.  The terms containing $d\zeta_R$ tend
to zero, including the cross term by Cauchy--Schwarz.  Dominated
convergence therefore yields
\[
 g\|\psi w_-\|_2^2\le C\delta^2\|\psi w_-\|_2^2,
\]
so $w_-=0$.  On an open subdomain the same compact-test argument
applies when the negative part has zero boundary trace, using the
Dirichlet form gap.  In particular it applies to nonnegative boundary
values.  Apply comparison to both signs of the difference of two
bounded solutions to obtain uniqueness.  If $M=1+\|W\|_\infty$, then
$L(s-M^{-1})=1-W/M\ge0$, and comparison gives $s\ge M^{-1}$.

Euclidean products retain bounded geometry and polynomial volume,
and their form gap follows by slicing.  The preceding proof therefore
applies verbatim; uniqueness identifies their solution with the
pullback of $s$.
\end{proof}

\begin{lemma}[Exact solutions and compatible approximations]\label{regular:solutions}
Fix \emph{all} the corrections of \Cref{regular:correction} and let
$s_A$ denote the exact bounded solution $L_As_A=1$.  There exist
positive smooth compatible functions $v_A$ with
\begin{equation}\label{regular:strict-supersolution}
 L_Av_A\ge\tfrac12
\end{equation}
on every packet.  They have positive infimum and bounded derivatives
packetwise.  These $v_A$ are approximations; they are not relabeled as
the original exact solutions $s_A$.
\end{lemma}

\begin{proof}
We first establish a uniform end estimate over a whole face, which
also avoids any compactness ambiguity at its deeper ends.  On an
exact elementary collar $B\times[T,\infty)$, the difference
$w=s_A-s_B$ is bounded and satisfies $L_Aw=0$.  Set
$m_B=\inf s_B>0$ and choose $a>0$ with
$Ca^2\|s_B\|_\infty\le1/2$.  The function
\[
 b(x,t)=e^{-a(t-T)}s_B(x)
\]
satisfies $L_Ab=e^{-a(t-T)}(1-Ca^2s_B)\ge\tfrac12e^{-a(t-T)}$.
Comparison on the half collar, using its sliced Dirichlet gap and
the Dirichlet comparison argument in the proof of
\Cref{regular:resolvent}, yields
\begin{equation}\label{regular:end-decay}
 |s_A-s_B|(x,t)\le
 \frac{\|s_A-s_B\|_\infty}{m_B}\,
 e^{-a(t-T)}s_B(x).
\end{equation}
The boundary inequalities hold at $t=T$ for both signs.  Uniform
interior elliptic estimates for the homogeneous difference give the
same exponential decay in $C^2$ on $t\ge T+1$.  All constants are
uniform in $x\in B$, including its deeper ends.  Thus the estimate holds
on each elementary collar, also where collars overlap.  Exact agreement
of the approximations on those overlaps will come from the compatible
face copies used in the following construction.

All operator coefficients and every exact $s_A$ are now fixed.  Let
$M_A$ be a finite bound for $L_A:C^2\to C^0$ in uniform charts, and
choose a target
\[
 \tau_A<\min\{\tfrac14\inf s_A,\ (4M_A)^{-1}\}.
\]
The finite family of collar blending cutoffs has a $C^2$ operator
bound $H_A\ge1$, independent of translation of their starts.  Include in $H_A$ the fixed finite factor for the number of copies.
For each packet $B$, let $\varepsilon_B$ be the minimum of
\[
 \frac{\tau_A}{\prod_{i=0}^{\ell-1}(4H_{A_i})}
 \qquad
 (A=A_0>A_1>\cdots>A_\ell=B),
\]
over all directed coface-to-face chains ending at $B$, including the
length-zero chain. There are finitely many such chains, so this
minimum is positive. Appending the edge $A>B$ to a minimizing chain
for $A$ gives $\varepsilon_B\le\varepsilon_A/(4H_A)$.
All these numbers are fixed after the corrections and operator norms,
and before constructing any $v_A$.

Construct $v_A$ from deeper faces upward.  Keep $s_A$ in the central
part and use the ordinary weights in the proof of
\Cref{regular:partition} to blend to the already constructed copies
of $v_B$.  Move the starts until
$\|s_A-s_B\|_{C^2}<\varepsilon_A/(4H_A)$ on all participating
collars, using \Cref{regular:end-decay}.  The persistent difference
$v_B-s_B$ is paid by its previously allocated tolerance, not by moving
the start.  Start also beyond the finitely many compatibility
thresholds for the $v_B$'s, so their copies agree on overlaps.
The derivative bounds for the weights then give
\[
 \|v_A-s_A\|_{C^2}\le\varepsilon_A\le\tau_A.
\]
The result is an exact face copy sufficiently far out, is smooth,
and has bounded higher derivatives by compact parts and finitely
many smooth products.  Finally $v_A\ge\tfrac34\inf s_A>0$ and
$L_Av_A\ge1-M_A\tau_A\ge3/4$, which implies
\Cref{regular:strict-supersolution}.  No geometry radius, correction,
or operator is changed during this approximation.
\end{proof}

\subsection{Restoring the invariant}

The passage from a positive solution to scalar-curvature control by a
circle warp appears in
\citet[Theorem~4\textup{(ii)}]{fischerColbrieSchoen1980}.
Here bounded positive supersolutions must also be exactly compatible
with every packet end; \Cref{regular:solutions} supplies these functions.

\begin{proposition}[The full invariant step]\label{regular:step}
Assume the stage $r$ invariant and the conclusion of
\Cref{shield:global}.  For every nonempty parent packet of base
dimension at least two, processing $j_0$ produces the compatible
smooth split packets of
\Cref{packet:split}, with the required forced signs and continuous
projection maps.  There are smooth compatible old metrics and $399$
new torus factors on every packet such that the profiles, cores,
potentials, and scalar curvature satisfy the stage $r+1$ invariant.
This includes a possible one-dimensional output, on which the scalar
lower bound is global.  After the radial split the packets have the
complete compact-part/product-end geometry of
\Cref{packet:geometry}.
\end{proposition}

\begin{proof}
Use \Cref{regular:partition,regular:seam,regular:profiles} for the
smooth base metrics and the profile data, and
\Cref{regular:correction,regular:solutions} for $v_A$.  Append
$k=399$ circle variables, each with length function
$v_A^{2/(k+1)}$.  The warped-product identity is
\[
 R_{\rm new}=R_{\rm old}-\frac{4k}{k+1}
                    \frac{\lap_{\rm old}v_A}{v_A}.
\]
Indeed a common scale $f$ contributes
$-2k f^{-1}\lap_{\rm old}f-k(k-1)f^{-2}|df|^2$; substitution of
$f=v_A^{2/(k+1)}$ cancels the separate $|dv_A|^2$ term.
The old Laplacian on invariant functions is the weighted base
Laplacian, including the product of all old warp lengths.  Here $k$ counts only the newly appended circle factors. Any unit
factors used earlier to compare old lists on a finite diagram have
zero logarithmic derivative and leave both $\lap_{\rm old}$ and this
coefficient unchanged. Since $4k/(k+1)=C$, \Cref{regular:Vstar,regular:strict-supersolution} give
\[
 R_{\rm new}=V_A^*-K_A+\frac{L_Av_A}{v_A}
 \ge V_A^*\quad\hbox{on }D_A.
\]
Here $K_A=0$ on a neighborhood of $D_A$.  Adding these factors does
not change base distances or profile gradients.  The new factors
are positive, bounded above and away from zero, and exactly compatible.
The same number is appended in every status at the stage, so all
end identities compare the same list of factors.  This proves the
analytic invariant, including when the new required threshold is zero.
The componentwise metric operations preserve every unused old circle
length identically equal to $1$: restriction, product extension, and
$\beta_s\mapsto\beta_{\phi(x,s)}$ all fix that constant entry.  The
$399$ appended lengths are the only common new entries.  Likewise
restriction and cylinder pullback preserve a constant pole profile and
its full core, and \Cref{regular:constant-preservation} preserves the
specified multiplied constant potentials.

The zero wall lies in $\operatorname{int}G_{j_0}$ and all current
forced signs hold by \Cref{regular:partition}; earlier forced cores
are carried by the continuous projection and retain their pole
neighborhoods.  Thus this is exactly the regular compatible splitting
operation required by the packet construction.  At the radial split,
the preserved alternating compact barrier bands separate all wall
components into bounded transition shells and trap every closed
sign-side component between two separating bands (with a compact
initial block).  These components are compact before prolongation.
The normal collars and smooth metrics used above are products in all
old ends.  \Cref{packet:split,packet:geometry} therefore supply the
claimed geometry after prolongation.  The main argument applies the
one-dimensional obstruction and reconstructs the compatible partitions in the original manifold.
\end{proof}
